\section{Free spaces and the graph-space obstruction}
\label{sec:free-graph}

The metric construction will provide a bounded linear operator
\(q:E\to U\) that is completely continuous, with \(E\) weakly
sequentially complete. This section explains why those two properties
exclude \(c_0(\ell_2)\) from the associated graph space. We first
record the free-space facts that produce \(E\), then prove the
obstruction without assuming compactness of any operator into
\(c_0(U)\).

\subsection{Free-space facts and the two external inputs}

For a pointed metric space \((M,d,0)\), let \(\Lip_0(M)\) be the
Banach space of real Lipschitz functions \(f:M\to\R\) with \(f(0)=0\),
normed by \(\Lip(f)\). Let \(\delta(x)\in\Lip_0(M)^*\) be evaluation
at \(x\), and define
\[
 \F(M)=\overline{\spanlin}\{\delta(x):x\in M\}
       \subset\Lip_0(M)^*.
\]
The evaluation and molecule construction goes back to Arens and Eells
\cite[Sections~2--3]{AE56}; the notation here is the later
Lipschitz-free formulation. In particular, \(\delta(0)=0\).
When a Lipschitz function \(f\) need not vanish at \(0\), the pairing
\(\langle f,\mu\rangle\) means \(\langle f-f(0),\mu\rangle\).

The evaluation map preserves distances:
\begin{equation}
\label{eq:delta-isometry}
 \|\delta(x)-\delta(y)\|=d(x,y).
\end{equation}
The Lipschitz inequality gives the upper bound. For the reverse bound,
test against \(t\mapsto d(t,y)-d(0,y)\), which has Lipschitz constant
at most one. If \(M\) is separable, choose a countable dense set
\(D_0\subset M\). Equation~\eqref{eq:delta-isometry} shows that finite
rational linear combinations of \(\delta(D_0)\) are dense in
\(\F(M)\); hence \(\F(M)\) is separable.

If \(K\subset M\) contains \(0\), the closed span of \(\delta(K)\)
in \(\F(M)\) is isometric to \(\F(K)\). To see this, a real
\(L\)-Lipschitz function \(f\) on \(K\) extends with the same
Lipschitz constant by the McShane formula
\cite[Lemma~3.1 and the following paragraph]{Kalton04}
\[
 \widetilde f(x)=\inf_{z\in K}\bigl(f(z)+L\,d(x,z)\bigr).
\]
For \(L=0\) use the constant extension. For \(L>0\), the Lipschitz
inequality on \(K\) makes the displayed infimum equal to \(f\) there,
and the triangle inequality gives the \(L\)-Lipschitz bound on \(M\).
The infimum is finite because \(0\in K\) supplies an upper bound and
the triangle inequality supplies a lower bound. Thus the suprema
defining the two free-space norms agree on finite evaluation sums
supported in \(K\), and then on their closures. In particular,
\begin{equation}
\label{eq:compact-pairing}
 |\langle f,\nu\rangle|
 \le \Lip(f|_K)\|\nu\| \qquad(\nu\in\F(K)).
\end{equation}

For any Banach space \(H\) and any Lipschitz map \(D:M\to H\) with
\(D(0)=0\), there is a unique bounded linear map
\begin{equation}
\label{eq:linearization}
 \widehat D:\F(M)\to H,\qquad
 \widehat D\delta(x)=D(x),\qquad
 \|\widehat D\|=\Lip(D).
\end{equation}
Indeed, for a finite evaluation sum, define
\(\widehat D(\sum_j a_j\delta(x_j))=\sum_j a_jD(x_j)\).
Testing the latter vector against a norm-one functional in \(H^*\)
and using the free-space norm gives
\[
 \left\|\sum_j a_jD(x_j)\right\|
 \le \Lip(D)\left\|\sum_j a_j\delta(x_j)\right\|.
\]
This inequality makes the definition independent of the chosen
representation and extends it continuously to \(\F(M)\).
Testing the extension on \(\delta(x)-\delta(y)\) proves the reverse
norm bound. This metric-domain linearization is recorded in
\cite[Lemma~3.2]{Kalton04}; for Banach-space domains see also
Godefroy and Kalton \cite[Lemma~2.5]{GK03}.

A bounded linear operator is \emph{completely continuous} if it sends
each weakly null sequence to a norm-null sequence. A Banach space is
\emph{weakly sequentially complete} if each weakly Cauchy sequence
has a weak limit in the space. We use the following two external
results with exactly the displayed hypotheses.

\begin{theorem}[Compact reduction for a weakly null sequence]
\label{thm:compact-reduction}
Let \((M,d,0)\) be a complete pointed metric space, and let
\(\mu_i\to0\) weakly in \(\F(M)\). For every \(\tau>0\), there is a
compact \(K\subset M\), with \(0\in K\), such that
\[
 \dist(\mu_i,\F(K))<\tau\qquad\hbox{for every }i.
\]
\end{theorem}

This is the sequential consequence of the tightness theorem of
Aliaga, No\^us, Petitjean and Proch\'azka
\cite[Theorem~2.3, arXiv version cited in the bibliography]{ANPP21}.
Their theorem states that every weakly precompact subset \(W\) of
\(\F(M)\), for complete \(M\), is tight: for each \(\epsilon>0\)
there is compact \(K_0\subset M\) with
\(W\subset\F(K_0)+\epsilon B_{\F(M)}\), where for an arbitrary
subset \(K_0\) this notation means its closed evaluation span inside
\(\F(M)\).
Here weak precompactness means that every sequence in \(W\) has a
weakly Cauchy subsequence. The set \(\{0,\mu_1,\mu_2,\ldots\}\)
has that property: any sequence in it has either a constant
subsequence or a subsequence whose indices tend to infinity, which
converges weakly to zero. Apply the theorem with \(\epsilon=\tau/2\)
and enlarge \(K_0\) by \(0\). This gives the strict inequality above.
Kalton's finite-radius localization lemma
\cite[Lemma~4.5]{Kalton04} is a precursor to this complete-space
compact reduction.

\begin{theorem}[Weak sequential completeness of a free space]
\label{thm:free-wsc}
If \(M\) is a superreflexive real Banach space, then \(\F(M)\) is
weakly sequentially complete. In particular, this applies to every
real Hilbert space \(M\).
\end{theorem}

This is \cite[Corollary~2.5, arXiv version cited in the
bibliography]{ANPP21}. Its compact-set input is the Main Theorem of
Kochanek and Perneck\'a \cite{KP18}, which concerns free spaces over
compact subsets of superreflexive Banach spaces, not directly the
whole Hilbert space.

\subsection{Weighted graph spaces}

Let \(A,H\) be real Banach spaces and \(Q:A\to H\) bounded
and linear. Define
\begin{align}
\label{eq:graph-space}
 Z_Q=\bigl\{s=(s_i)_{i\ge1}:{}&
 s_i\in A,\quad \sum_{i\ge1}2^{-i}\|s_i\|<\infty,\quad
 (Qs_i)_{i\ge1}\in c_0(H)\bigr\},\\
\label{eq:graph-norm}
 \|s\|_{Z_Q}
 ={}&\max\left\{\sum_{i\ge1}2^{-i}\|s_i\|,
                         \sup_{i\ge1}\|Qs_i\|\right\}.
\end{align}
The membership condition is part of the definition; finiteness of the
displayed maximum alone would not require the \(Q\)-images to tend to
zero.

\begin{lemma}
\label[lemma]{lem:graph-banach}
\(Z_Q\) is a Banach space, and it is separable if \(A,H\) are
separable. The coordinate map
\(p_i:Z_Q\to A\), \(p_i(s)=s_i\), has norm at most \(2^i\), and
\[
 J_Q:Z_Q\to c_0(H),\qquad J_Q(s)=(Qs_i)_i,
\]
has norm at most one.
\end{lemma}

\begin{proof}
Let \(\ell_1(2^{-i};A)\) denote the weighted sum space with norm
\(\sum_i2^{-i}\|s_i\|\). The map \(s\mapsto(2^{-i}s_i)_i\)
is an isometry from this space onto \(\ell_1(A)\), so it is complete.
The map
\[
 s\longmapsto \bigl(s,(Qs_i)_i\bigr)
\]
identifies \(Z_Q\) isometrically with a subspace of
\(\ell_1(2^{-i};A)\oplus_\infty c_0(H)\). This subspace is closed.
In fact, suppose \((s^{(k)},(Qs_i^{(k)})_i)\) converges there to
\((s,y)\). For each fixed \(i\),
\[
 \|s_i^{(k)}-s_i\|
 \le 2^i\|s^{(k)}-s\|_{\ell_1(2^{-j};A)}
 \longrightarrow0.
\]
Continuity of \(Q\) gives \(Qs_i^{(k)}\to Qs_i\), whereas convergence
in \(c_0(H)\) gives \(Qs_i^{(k)}\to y_i\). Thus \(y_i=Qs_i\) for all
\(i\), as required. Completeness follows.

If \(A,H\) are separable, truncating a sequence in \(Z_Q\) gives
convergence in both parts of
\eqref{eq:graph-norm}: the weighted sum has a vanishing tail, and the
\(Q\)-image is norm-null. Choose a fixed countable dense subset of
\(A\). The finitely supported sequences with entries in that subset
form a countable dense subset of \(Z_Q\), because each of their
finitely many entries may be approximated in \(A\) and \(Q\) is
continuous. Finally, the two claimed operator bounds are immediate
from \eqref{eq:graph-norm}.
\end{proof}

The next lemma records exactly the consequence of weak sequential
completeness that the graph obstruction needs. It does not assert
that every operator into a \(c_0\)-sum is compact.

\begin{lemma}
\label[lemma]{lem:no-c0-sum}
If \(E\) is a weakly sequentially complete real Banach space, then
\(\ell_1(E)\) contains no closed subspace linearly isomorphic to
ordinary \(c_0\).
\end{lemma}

\begin{proof}
Every finite direct sum of \(E\) is weakly sequentially complete:
the coordinates of a weakly Cauchy sequence are weakly Cauchy in
\(E\), hence have weak limits, and the finite list of limits is the
weak limit in the sum. The property also passes to a closed linear
subspace \(F\) of a weakly sequentially complete space \(G\).
Indeed, a weakly Cauchy sequence in \(F\) is weakly Cauchy in \(G\);
its weak limit lies in \(F\), since a closed linear subspace is
weakly closed by Hahn--Banach separation. Every functional on \(F\)
extends to \(G\) by Hahn--Banach, so the same limit is weak in \(F\).

Ordinary \(c_0\) is not weakly sequentially complete. If
\((e_k)\) is its usual unit basis, the partial sums
\(e_1+\cdots+e_n\) are weakly Cauchy: for every
\((a_k)\in c_0^*=\ell_1\), their pairings are the Cauchy partial sums
\(\sum_{k\le n}a_k\). A weak limit would have every coordinate equal
to one, which is not an element of \(c_0\).

Suppose now that \(A:c_0\to\ell_1(E)\) is a linear embedding with
\(\|Ax\|\ge b\|x\|\), where \(b>0\). Let \(P_N\) be projection onto
the first \(N\) coordinates, with \(P_0=0\). No \(P_NA\) is bounded
below, because otherwise its closed range would be a copy of \(c_0\)
in the weakly sequentially complete finite sum \(E^N\).

Starting with \(N_0=0\), choose unit vectors \(x_j\in c_0\) and
integers \(N_j>N_{j-1}\) recursively so that
\[
 \|P_{N_{j-1}}Ax_j\|<b/8,\qquad
 \|(I-P_{N_j})Ax_j\|<b/8.
\]
The first inequality is possible because \(P_{N_{j-1}}A\) is not
bounded below; after choosing \(x_j\), its \(\ell_1(E)\) tail gives
the second. Put \(z_j=(P_{N_j}-P_{N_{j-1}})Ax_j\). Then
\[
 \|Ax_j-z_j\|<b/4,\qquad \|z_j\|>3b/4.
\]
The \(z_j\) have disjoint coordinate supports. For any finite scalar
family \((\alpha_j)\),
\begin{align*}
 \left\|\sum_j\alpha_jAx_j\right\|
 &\ge \left\|\sum_j\alpha_jz_j\right\|
       -\sum_j|\alpha_j|\|Ax_j-z_j\|\\
 &=\sum_j|\alpha_j|\|z_j\|
       -\sum_j|\alpha_j|\|Ax_j-z_j\|
 \ge \frac b2\sum_j|\alpha_j|.
\end{align*}
Consequently,
\[
 \frac{b}{2\|A\|}\sum_j|\alpha_j|
 \le\left\|\sum_j\alpha_jx_j\right\|_{c_0}
 \le\sum_j|\alpha_j|.
\]
The finite-sum map therefore extends by completeness to a bounded
linear embedding of \(\ell_1\) into \(c_0\), with closed range.
This is impossible. If \(F\) is any closed subspace of \(c_0\), the
restriction map \(c_0^*=\ell_1\to F^*\) is onto by Hahn--Banach,
so \(F^*\) is separable. A subspace isomorphic to \(\ell_1\) would
instead have dual isomorphic to \(\ell_\infty\), which is
nonseparable: the uncountable set of all \(0\)-\(1\) sequences is
pairwise separated by distance one. This contradiction proves the
lemma.
\end{proof}

\begin{lemma}
\label[lemma]{lem:graph-obstruction}
Let \(E,U\) be real Banach spaces, with \(E\) weakly sequentially
complete, and let \(q:E\to U\) be a completely continuous bounded
linear operator.
Then \(Z_q\) contains no closed subspace linearly isomorphic to
\(c_0(\ell_2)\).
\end{lemma}

\begin{proof}
Suppose that \(S:c_0(\ell_2)\to Z_q\) is a linear embedding with
\(\|Sx\|_{Z_q}\ge b_0\|x\|\), \(b_0>0\).
We will select a scalar copy of \(c_0\) on which the \(c_0(U)\)
part of \(S\) has norm at most \(b_0/2\). The weighted sum part of
the graph norm must then retain the lower bound \(b_0\), giving a
copy of \(c_0\) in \(\ell_1(E)\), contrary to
\Cref{lem:no-c0-sum}.

Write
\(I_n:\ell_2\to c_0(\ell_2)\) for injection into the \(n\)th
coordinate and put
\[
 T_n=J_qSI_n:\ell_2\to c_0(U).
\]
We first prove that each \(T_n\) fails to be bounded below.

Fix \(n\) and an orthonormal sequence \((e_k)\) in \(\ell_2\).
It is weakly null, since the scalar coefficients of a fixed Hilbert
vector against an orthonormal sequence are square summable. For each
fixed \(i\), boundedness of \(p_iSI_n\) therefore gives
\(p_iSI_ne_k\to0\) weakly in \(E\). Complete continuity of \(q\)
implies
\[
 \|(T_ne_k)_i\|=\|q\,p_iSI_ne_k\|\longrightarrow0
 \qquad\hbox{for each fixed }i.
\]
Set \(y_k=T_ne_k\), so \(\|y_k\|\le\|T_n\|\).
Choose \(k_j\) increasing and integers
\(0=m_0<m_1<m_2<\cdots\) as follows. Having chosen \(m_{j-1}\),
choose \(k_j\) so large that
\[
 \max_{i\le m_{j-1}}\|(y_{k_j})_i\|<2^{-j}.
\]
The maximum over no coordinates is zero. Then choose \(m_j>m_{j-1}\)
so large that \(\sup_{i>m_j}\|(y_{k_j})_i\|<2^{-j}\), using
\(y_{k_j}\in c_0(U)\). Let \(z_j\) be the restriction of \(y_{k_j}\)
to coordinates \(m_{j-1}<i\le m_j\). Thus
\[
 \|y_{k_j}-z_j\|<2^{-j},\qquad \|z_j\|\le\|T_n\|.
\]
Because the \(z_j\) have disjoint coordinate supports, for all \(N\)
\[
 \left\|\sum_{j=1}^N T_ne_{k_j}\right\|
 \le \left\|\sum_{j=1}^N z_j\right\|
          +\sum_{j=1}^N2^{-j}
 \le \|T_n\|+1.
\]
The corresponding input sum has Hilbert norm \(\sqrt N\).
No positive lower bound for \(T_n\) is compatible with these two
estimates.

Choose a unit vector \(t_n\in\ell_2\) such that
\(\|T_nt_n\|\le b_0\,2^{-n-1}\). The map
\[
 R:c_0\to c_0(\ell_2),\qquad R(a)=(a_nt_n)_n
\]
is a linear isometry. For finitely supported \(a\),
\[
 \|J_qSRa\|
 \le\sum_n |a_n|\|T_nt_n\|
 \le\frac{b_0}{2}\|a\|_\infty.
\]
Truncation in \(c_0\) and continuity extend this estimate to every
\(a\in c_0\). On the other hand
\(\|SRa\|_{Z_q}\ge b_0\|a\|_\infty\).
The second part of the maximum in \eqref{eq:graph-norm} is at most
\((b_0/2)\|a\|_\infty\), so the first part must satisfy
\[
 \sum_i2^{-i}\|(SRa)_i\|\ge b_0\|a\|_\infty.
\]
For \(a=0\) this is immediate; for \(a\ne0\) it follows because the
second part is strictly below the required lower bound for the
maximum. Hence
\[
 a\longmapsto \bigl(2^{-i}(SRa)_i\bigr)_i
\]
is a bounded linear embedding of \(c_0\) into \(\ell_1(E)\).
Boundedness comes from \(\sum_i2^{-i}\|(SRa)_i\|\le\|SRa\|_{Z_q}\),
and the displayed lower bound makes the range closed. This
contradicts \Cref{lem:no-c0-sum}.
\end{proof}
